% ===== §6 Non-Fullness and Renewal =====
\section{Non-Fullness and Renewal}
\label{p20:sec:nonfullness}

\noindent
The third water-virtue is the most quietly demanding, because it asks something of a relation precisely when nothing is wrong with it. ``To hold a thing brim-full is not so good as to stop in time'' (ch.~9); ``it is just because it is never full that it can wear out and be renewed'' (ch.~15). The counsel is against completion --- against the filling of a relation to the point where it is finished, settled, brought to a state from which the only motion is decline. Where non-contention concerns what one grasps and the low place concerns where one attends, non-fullness concerns when one \emph{stops} --- and the claim is that a relation kept deliberately unfilled, refused the satisfaction of being complete, is the one that can keep renewing itself, while the relation filled to the brim has nowhere to go but to spill or stagnate.

\subsection{The ethics of the unfilled}
\label{p20:sub:nonfull-ethics}

Stated as an ethic, non-fullness is the refusal of two closures that masquerade as goods. The first is the closure of self-satisfaction: the relation that takes itself to have arrived, to have become what it was for, and so stops generating because it believes the work done. The second is the closure of completion-over-the-other: the taking of the other as fully known, read to the bottom, settled into a type. Against both, non-fullness keeps a deliberate remainder --- leaves room, declines to fill the last of the space, holds the relation open at the very point where it could be sealed. There is a thin secular wisdom that has always known this. Socratic knowledge is knowledge of one's ignorance, a refusal of the fullness of supposed completeness that keeps inquiry alive; the moment one is full of knowing, one has stopped learning, and the same holds of a person as of a subject. Game theory supplies an unexpectedly exact corollary on the side of conflict: the strategy that drives its counterpart to the wall, that fills its advantage to the brim and leaves the other no room, destroys the repeated relation it depends on; the grim and total victory is the one that ends the game. The converse is the water-virtue's counsel --- do not fill the advantage to the brim, do not press the other to the wall, leave the room that lets the relation continue. Non-fullness is, in this register, simply the recognition that in a relation meant to last, the unused remainder is not waste but the condition of continuance.

\subsection{Non-fullness in intimacy}
\label{p20:sub:nonfull-intimacy}

In intimacy the unfilled remainder has a precise and tender content: it is the room left for the other to be more, and other, than one has so far known. The temptation non-fullness resists is the temptation to complete one's knowledge of the partner --- to settle who they are, to fix them in a character one has inferred, to meet today the person one met yesterday. This temptation wears the mask of intimacy itself; ``I know you better than anyone'' is offered, and often received, as the height of closeness. But knowledge held as complete is knowledge that has stopped, and a partner met as fully known is a partner denied the room to change. This series has argued at length that the other is, in principle, never exhausted --- that happiness shared is hermeneutic, given as the continual interpretation of a meaning-world that does not bottom out, and that the resonance of two such worlds depends on each remaining inexhaustible to the other rather than being annexed and used up \parencite{paperxix}. Non-fullness is the ethical form of that ontological fact. To keep the relation unfilled is to keep meeting the partner as still-unfinished, to hold one's knowledge of them open to revision, to leave room for the self they have not yet become and the depth one has not yet reached. ``Worn out and renewed'' is, in intimacy, exactly this: a relation that does not seal its account of the other is one that can keep discovering them, and a relation that can keep discovering is one that does not exhaust its own generativity. The unfilled remainder is where tomorrow's intimacy comes from; to fill it --- to know the other completely, to settle the relation entirely --- is to spend the future for the sake of a present completeness, and to be left, once the knowing is complete, with nothing further to generate.

\subsection{Non-fullness as the avoidance of solidification}
\label{p20:sub:nonfull-formal}

The formal companion gives this its exact dynamical form, and the form is worth stating because it shows non-fullness to be not a mood but a structural condition. In the companion's dynamical layer, a generative relation is a circulating flow, and there is a definite way for such a flow to die: when the parameter governing whether the relation's accumulated history is forgotten or compounded tips toward compounding, a positive feedback locks the flow into a fixed channel, its circulation falls to zero, and the living, surplus-generating cycle collapses into a frozen, merely conservative one --- a structure that still exists but no longer generates. The companion calls this solidification, and identifies it as the dynamical face of the bad cycle \parencite{wan2026braid}. Non-fullness is the ethical name for the maintenance of a relation away from that frozen fixed point. To fill a relation to completion --- to settle it, to let its accumulated history harden into a fixed channel through which everything must now run --- is solidification described from inside; the relation that ``knows'' itself and the other completely is the relation whose circulation has gone to zero, that no longer generates because it has frozen into what it already is. ``Worn out and renewed'' is the refusal of that freezing: the continual reopening of a flow that would otherwise set, the deliberate preservation of the remainder that keeps the channel from hardening. This subsumes, and makes precise, the criterion of the good cycle developed earlier in this series, where the good relation was distinguished by its structural refusal of closure --- its accumulation of positive surplus along a path that does not return to its origin \parencite{paperix}. Non-fullness is that refusal of closure stated as a virtue of conduct: the good cycle is the one that does not fill itself shut, and the water-virtue of non-fullness is the discipline of keeping it open. The full dynamical statement is given in the translation layer (\xref{p20:sub:formal-solidification}); here it is enough that the counsel against brimming-full and the dynamics of solidification are one thing seen from two sides --- the ethics of the unfilled and the physics of the flow that does not freeze.
